% VAFT Beamer theme.
%
% This is the shared preamble the five hand-written decks in tutorial/ each
% carry a duplicated copy of. Centralising it here is the point of the QMD
% pipeline: one place defines VAFT's slide identity for every deck and both
% backends.
%
% Keep this in step with vaft.scss, which does the same job for the Reveal.js
% output. The two are the same design expressed twice; the sizes below are the
% point-equivalents of the Reveal defaults, so a slide looks the same in either.

% --- Typeface --------------------------------------------------------------
%
% Beamer's default is Latin Modern Sans, which has a noticeably smaller x-height
% than the Source Sans Pro that Reveal.js ships. Setting both sides to a
% Helvetica-metric face is what makes the two outputs actually look like the
% same deck rather than merely sharing a colour.
%
% helvet (Nimbus Sans) is in texlive-fonts-recommended, which CI already
% installs. Matching Reveal's bundled Source Sans Pro exactly would mean
% sourcesanspro from texlive-fonts-extra -- about 1.2 GB on a job that exists to
% be fast -- so both sides meet on Helvetica instead. See vaft.scss.
\usepackage[scaled=0.95]{helvet}
\renewcommand{\familydefault}{\sfdefault}

% --- Colour ----------------------------------------------------------------
% Matches $vaft-blue in vaft.scss and \definecolor in the .tex decks.
\definecolor{vaftblue}{RGB}{25,76,127}
\setbeamercolor{structure}{fg=vaftblue}
\setbeamercolor{title}{fg=vaftblue}
\setbeamercolor{frametitle}{fg=vaftblue}

% --- Type sizes ------------------------------------------------------------
%
% A 16:9 Beamer canvas is 90 mm tall; a Reveal slide is 700 logical px. One px
% is therefore 90/700/0.35278 = 0.3645 pt, and these are the Reveal defaults
% converted through it:
%
%   title 100px -> 36pt   frametitle 64px -> 23pt
%   subtitle    40px -> 15pt   body       40px -> 15pt
%
% Headings are bold on both sides: vaft.scss sets
% $presentation-heading-font-weight: 600, and Beamer's default series is
% regular, which was the most visible remaining difference between the outputs.
\setbeamerfont{title}{size=\fontsize{36}{42}\selectfont, series=\bfseries}
\setbeamerfont{subtitle}{size=\fontsize{15}{19}\selectfont, series=\mdseries}
\setbeamerfont{frametitle}{size=\fontsize{23}{27}\selectfont, series=\bfseries}
\setbeamerfont{normal text}{size=\fontsize{15}{19}\selectfont}

% The subtitle is subordinate to the title on both sides, not a second heading.
% Matches rgba(26, 26, 26, 0.75) in vaft.scss.
\definecolor{vaftsubtitle}{RGB}{70,70,70}
\setbeamercolor{subtitle}{fg=vaftsubtitle}

% Reveal.js draws round bullets; Beamer's default is a triangle. Matching
% them removes the last obvious tell that these are two different decks.
\setbeamertemplate{itemize item}{\textbullet}
\setbeamertemplate{itemize subitem}{--}

\setbeamertemplate{navigation symbols}{}
\setbeamertemplate{footline}[frame number]

% Section slides, which the hand-written decks had no mechanism for.
\setbeamertemplate{section page}{%
  \begin{centering}
    \usebeamerfont{title}\usebeamercolor[fg]{structure}\insertsectionhead\par
  \end{centering}%
}
\AtBeginSection[]{\frame{\sectionpage}}

% --- Speaker notes ---------------------------------------------------------
%
% Pandoc turns every ::: {.notes} block into \note{...}, and Beamer's default is
% to hide those. That is what the archival PDF wants. The presenter build passes
% \setbeameroption{show notes} instead, so the same notes become facing pages --
% authored once, in one source. See tutorial/presentations/README.md.
